\documentclass[11pt]{article}
\usepackage[margin=1in]{geometry}
\usepackage[T1]{fontenc}
\usepackage{lmodern,amsmath,amssymb,amsthm,mathtools,hyperref}
\setlength{\emergencystretch}{3em}
\newtheorem{theorem}{Theorem}[section]
\newtheorem{lemma}[theorem]{Lemma}
\newtheorem{corollary}[theorem]{Corollary}
\newtheorem{proposition}[theorem]{Proposition}
\newcommand{\E}{\mathbb E}
\newcommand{\tr}{\operatorname{tr}}
\title{Compatible modular codes, exact lifting,\\and concentration barriers}
\author{Joshua Mark\thanks{ORCID: \href{https://orcid.org/0009-0003-5905-2382}{0009-0003-5905-2382}. AI tools assisted with derivation, independent review, computation and exposition.}}
\date{20 November 2026 --- research draft; priority not established}
\begin{document}\maketitle

\begin{abstract}
We derive a parity rigidity theorem for modularly orthogonal sign
families of length $4q$, $q$ odd. In a self-orthogonal residue code
containing five real orthogonal rows, all sign words are mutually
compatible. Over a prime field an exact extension gives a self-dual
neighbor preserving all old sign words; at the known order twenty,
one such change raises the sign count from eight to twenty.
A second order-twenty example is terminal: no new sign row preserves
eleven of its twelve rows. An exact two-row replacement unlocks ten
new rows and a complete matrix, again in one self-dual neighbor.
An exact neighbor-incidence identity gives a quantitative surplus
criterion; in a codimension-one retained space the entire modular
sign pool splits between just two self-dual codes.
We also classify modular false positives at orders congruent to four
modulo eight, give exact monotone concentration moves and limits of
specified-target bridge barriers, and strengthen a low-degree model
obstruction to an explicit moment-error floor.
None of these results proves a sufficiently large sign count at
unknown orders. The full conjecture and bibliographic novelty remain
unresolved.
\end{abstract}
The standing distribution target remains exact unbiased pairwise
independence on $N-4$ bits and collision at least
$1/N-1/(1024N^4)$; the proven extraction implication is in
\path{../round6/collision_stability.tex}. Here we investigate an
alternative global residue-code model as well as its obstructions.
\section{Compatible codewords and a change of self-dual code}
\label{sec:compatible-codes}

Put $N=4q$ with $q\ge3$ odd. A sign class means a pair $\{v,-v\}$;
use the representative whose first coordinate is positive. Set
$\pi(v)=\prod_jv_j$, which is well-defined on classes because $N$ is even.
We call a family modularly orthogonal if all its pairwise inner products
are divisible by $q$. The following proof requires no primality.

\begin{theorem}[Two-parity rigidity]\label{code:parity}
In a modularly orthogonal family of distinct sign classes, the
classes of each fixed parity are pairwise orthogonal over the reals.
If both parities occur, each parity has at most four classes.
Consequently any such family with at least nine classes is a real
partial Hadamard matrix, and every family has at most $N$ classes.
The mixed-parity bound of eight is sharp for every odd $q\ge3$.
\end{theorem}
\begin{proof}
For distinct nonantipodal signs $u,v$, their inner product is even,
divisible by odd $q$, and strictly between $-4q$ and $4q$. It is
therefore $0$ or $\pm2q$. Since
\[
 \pi(u)\pi(v)=(-1)^{(N-u\cdot v)/2},
\]
equal parities force zero and opposite parities force $\pm2q$.
Within one parity the normalized rows are orthonormal. Against
any normalized row of the opposite parity, every squared inner
product is $1/4$. Bessel's inequality limits the first class to
four rows, and then does the same for the second.
If only one parity occurs there are at most $N$ orthogonal rows;
otherwise there are at most eight, which is less than $N$.

For sharpness start with the four rows of $H_4$. Adjoin the four
antipodal classes of sign vectors of the form
$\frac12\sum_{i=1}^4\sigma_i h_i$ that have all entries $\pm1$.
The two sets are orthogonal bases of $\mathbb R^4$ whose cross
inner products are $\pm2$. Repeat each coordinate $q$ times.
This gives four classes of each parity, with cross products
$\pm2q$ and within-class products zero.
\end{proof}

This is an elementary rigidity result within the classical subject
of modular Hadamard matrices, not a priority claim; see Marrero and
Butson, \emph{Modular Hadamard matrices and related designs} (1973),
\url{https://doi.org/10.1016/0097-3165(73)90073-3}.
The modulus here grows with the target order. Existence theorems
for a fixed small modulus cannot simply supply this hypothesis.

\begin{corollary}[Five rows remove compatibility tests]\label{code:five}
Let $C$ be a submodule of $(\mathbb Z/q\mathbb Z)^N$ with
$C\subseteq C^\perp$, and let
\[
 S(C)=\{v\in\{\pm1\}^N:v\bmod q\in C,\ v_1=1\}.
\]
If $C$ contains five mutually orthogonal sign rows, every pair of
distinct members of $S(C)$ is orthogonal over the reals.
Thus $5\le |S(C)|\le N$.
\end{corollary}
\begin{proof}
The five prescribed orthogonal classes all have the same parity.
One opposite-parity class in $C$ would violate the four-row Bessel
bound. All sign classes in $C$ consequently have this same parity,
and Theorem~\ref{code:parity} applies.
\end{proof}

This changes the task from finding a clique among modular candidates
to counting sign words in one self-orthogonal code. It does not prove
that the count is large. In particular a dual kernel $P^\perp$ of a
short frame is usually not self-orthogonal; this corollary must not
be applied to that entire candidate set.
Five initial orthogonal rows do exist without assuming $H_N$:
write $N=12+8k$ and concatenate the first five rows of a fixed
$H_{12}$ with $k$ copies of the first five rows of $H_8$.
Their modular row span is self-orthogonal, but this construction
alone supplies no useful lower bound beyond five sign classes.

\subsection{A precise restricted-alphabet existence criterion}
The elementary three-column completion theorem proved in
\path{../round5/collision_entropy.tex} says that $N-3$ orthogonal
sign rows of length $N$ can be completed to a full Hadamard matrix.
For $N\ge12$, $N-3\ge9$. The preceding theorem therefore gives:
\[
 \begin{split}
 H_N\text{ exists}\quad\Longleftrightarrow\quad&
 \text{some self-orthogonal submodule }C\subseteq
 (\mathbb Z/q\mathbb Z)^N\\[-2pt]
 &\text{has }|S(C)|\ge N-3.
 \end{split}
\]
For the forward implication take the module generated by the rows
of $H_N$, whose inner products are zero modulo $q$, including the
diagonal $N$. For the reverse implication the sign words are a
partial Hadamard matrix by Theorem~\ref{code:parity}, and complete
its $N-3$ rows. This is an equivalent target, not a count estimate.
One may equivalently require the module to be self-dual. Indeed the
row module of a full $H_N$ is self-dual even for composite odd $q$:
if $Hz=qy$ for integer lifts, then $H^TH=4qI$ implies
$4z=H^Ty$. Since four is invertible modulo $q$, the residue of
$z$ belongs to the row module. Self-orthogonality gives the reverse
inclusion. This is the classical Hadamard-code construction, not
a new self-dual-code existence theorem; see Munemasa and Tamura,
\emph{The codes and the lattices of Hadamard matrices},
\url{https://doi.org/10.1016/j.ejc.2011.11.007}.

Only odd $q$ are needed for the full conjecture: write an arbitrary
multiple of four as $4q\,2^k$ with $q$ odd, and tensor an order-$4q$
matrix with a Sylvester matrix of order $2^k$. The case $q=1$
is already Sylvester. Neither this reduction nor the criterion
establishes existence at the remaining odd multipliers.

\subsection{A self-dual neighbor preserving existing sign rows}
Now assume $q=p$ is an odd prime and work over $\mathbb F_p$.
Let $C=C^\perp$ be self-dual, with at least five orthogonal sign
classes in $S(C)$. If a sign vector $v$ is orthogonal over the
integers to every member of $S(C)$ but does not belong to $C$, put
\[
 C'=(C\cap v^\perp)+\langle v\rangle.
\]
\begin{proposition}[Exact code pivot]\label{code:pivot}
The space $C'$ is self-dual, $\dim(C\cap C')=N/2-1$, and
$S(C)\subsetneq S(C')$. Conversely, if no real sign extension of
$S(C)$ exists, no self-orthogonal code preserving every member
of $S(C)$ can acquire a new sign class.
\end{proposition}
\begin{proof}
The functional $c\mapsto c\cdot v$ on $C$ is nonzero since
$C=C^\perp$ and $v\notin C$, so its kernel has dimension $N/2-1$.
Also $v\cdot v=N=0$ in $\mathbb F_p$. The displayed sum is
self-orthogonal of dimension $N/2$, hence self-dual.
It contains the old sign words because their real inner products
with $v$ vanish, and it contains the new class of $v$.
The assumption $v\notin C$ is automatic for a new real extension:
an extension is not one of the classes being extended, whereas
membership in $C$ would put its normalized representative in $S(C)$.
The converse follows from the same corollary: any acquired class
would be a real extension of all preserved classes.
\end{proof}

Thus repeated extension-guided pivots strictly increase the sign
count and stop after at most $N-|S(C)|$ moves, either at a full matrix
or at a genuinely terminal partial frame. This is not an algorithm
that escapes terminal frames. Moving to a different code can remove
an artificial restriction, but cannot make an incompatible new row
orthogonal to all existing rows.

\subsection{Exact known-order experiments}
The verifier exhausts all self-dual ternary codes containing the
first five rows of a specified Paley $H_{12}$. Their span has rank
five, and there are exactly two self-dual extensions, each with
twelve sign classes. These are results for that prescribed seed,
not a classification of all five-row frames.

At order twenty, a deterministic finite-field construction extends
the first five rows of a Paley $H_{20}$ to a self-dual code over
$\mathbb F_5$. Exhausting its $2^9$ normalized sign assignments in
systematic coordinates yields only eight classes. This disproves
the suggestion that any self-dual extension containing five rows
automatically gives a full matrix. The eight-row frame has rank
eight modulo five. Enumerating its modular kernel gives forty
classes: the eight existing rows and thirty-two exact extensions.
Using the first extension in a single pivot gives
\[
             |S(C)|=8\quad\longrightarrow\quad |S(C')|=20.
\]
All twenty rows are checked pairwise over the integers.
The finite-field bases and the extension witness are recorded in
\path{codes/checks.json}.

Finally self-duality alone need not give even one sign word.
Over $\mathbb F_3$, the length-four code generated by
$(1,0,1,1)$ and $(0,1,1,-1)$ is self-dual and has no word of full
support. Its threefold direct sum is a self-dual length-twelve code
with no sign words. A useful general lower bound must use more than
dimension and self-duality. No count near $N$ has been proved at
unknown orders.
\section{A terminal code trap with an optimal two-row escape}

This certificate concerns order twenty, where Hadamard matrices are
already known. It gives an exact trap for monotone row addition and
an exact minimal-loss escape in the self-dual-code construction.

\begin{theorem}[Finite certified example]
There is a self-dual ten-dimensional code $C\subseteq\mathbb F_5^{20}$
with exactly twelve antipodal classes of sign words. Representatives
with first coordinate $+1$ form a real orthogonal $12\times20$ sign
frame $P$, and span $C$ over $\mathbb F_5$.
No further sign row is orthogonal to every row of $P$.
Every new sign class is orthogonal to at most ten old rows, so no
one-row deletion admits even one genuinely new row.
Deleting the two old rows indexed $3,8$ (zero based) permits a full
Hadamard completion of the remaining ten rows. Its mod-five code
$C'$ is a self-dual neighbor: $\dim(C\cap C')=9$.
One explicit neighbor pivot takes the sign count from twelve to twenty.
Thus two is the exact minimum number of old rows lost in such an escape.
\end{theorem}

The complete basis and twelve sign representatives are stored in
\texttt{search.json}. The construction starts with the first five rows
of the Paley order-twenty matrix and extends their span over
$\mathbb F_5$ using deterministic seed $2026112000$.
The construction is only a way to find the certificate; correctness
is checked from its stored entries. The code basis is self-orthogonal
of rank ten. Enumerating its $2^9=512$ systematic sign assignments
finds exactly twelve normalized sign words, of mod-five rank ten;
their real Gram matrix is $20I_{12}$.
If a new sign row were exactly orthogonal to $P$, its residue would
belong to $C^\perp=C$, hence to the already exhausted twelve classes.
Every such class fails orthogonality to itself, proving terminality.

An independent binary enumeration checks all $2^{19}=524288$
normalized sign words using only integer Hamming distances. Excluding
the twelve existing words, its compatibility counts are
\[
\begin{array}{c|rrrrrrrrrrr}
\text{compatible}&0&1&2&3&4&5&6&7&8&9&10\\\hline
\text{new classes}&262840&6840&27092&48948&67344&58976&34920&12600&3992&616&108
\end{array}
\]
No new word is orthogonal to eleven or twelve old rows. The 108
best words group by their two nonorthogonal old rows. For the pair
$(3,8)$, all ten available words are mutually orthogonal; together
with the ten retained old words they give Gram matrix $20I_{20}$.
Their normalized negative-coordinate bit masks are
\[
7676,126672,307898,444300,581982,645928,683178,816740,855582,886242.
\]
Bit $j$ represents coordinate $j$, with bit one denoting a minus sign.
The retained old rows have modular rank nine, the completed code
rank ten, and the union of the old and new code bases rank eleven.
These prove the claimed intersection dimension. If $v$ is the first
new row, then
\[
 C'=(C\cap v^\perp)+\langle v\rangle,
\]
as independently checked by row reduction.

The scripts \texttt{search.py} and \texttt{verify.py} use only the
Python standard library and original project code from the adjacent
\texttt{codes/check.py}. The latter also writes
\texttt{verification.json} and \texttt{exchange.json}. No external
implementation is imported. This example disproves the strategy of
always retaining all previously selected rows; it does not assert
that two-row neighbor escapes suffice in general or yield unknown
Hadamard orders.

\begin{theorem}[Elementary tensor persistence]
If a sign frame $P$ admits no orthogonal sign row and $H$ is a square
Hadamard matrix of order $t$, then $P\otimes H$ also admits no
orthogonal sign row. In particular this certificate gives terminal
$12t\times20t$ frames whenever such an $H$ is supplied.
\end{theorem}
\begin{proof}
Reshape a hypothetical new sign vector into a $20\times t$ sign
matrix $X$. Orthogonality to $P\otimes H$ says $PXH^T=0$, under
the corresponding vectorization convention. Since $H$ is invertible,
$PX=0$. Each column of $X$ would therefore be a sign vector
orthogonal to $P$, a contradiction.
\end{proof}
\section{An exact local count identity and a two-code decomposition}
\label{sec:neighbor-counts}

The preceding pivot suggests counting over code neighbors, rather
than asking for an extension preserving every row. The following
identity makes that approach quantitative. Here $p$ is an odd prime,
$N=4p$, $m=N/2$, $C=C^\perp$, and $G\subsetneq C$ has dimension
$s=m-h$, $h\ge1$. Put
\[
 r_C=|S(C)|,\quad b_G=|S(G)|,\quad
 M_G=|\{v\in\{\pm1\}^N:v_1=1,\ v\bmod p\in G^\perp\}|.
\]
Let $\mathcal N_G(C)$ be the self-dual codes $C'\ne C$ with
$\dim(C\cap C')=m-1$ and $G\subseteq C'$.

\begin{theorem}[Exact restricted-neighbor incidence count]
\label{code:neighbor-average}
There are $(p^h-1)/(p-1)$ codes in $\mathcal N_G(C)$, and
\[
 \boxed{\quad
 \sum_{C'\in\mathcal N_G(C)}\bigl(|S(C')|-r_C\bigr)
 =M_G-r_C-(r_C-b_G)p^{h-1}.\quad}
\]
In particular, if the right-hand side is positive, some neighbor
strictly increases the sign count while preserving $G$.
\end{theorem}
\begin{proof}
A hyperplane $D$ of $C$ containing $G$ has $D^\perp/D$ a nondegenerate
two-dimensional hyperbolic space. It has exactly two isotropic lines,
one giving $C$ and the other a unique neighbor $C'$.
To see the last assertion directly, choose isotropic $e$ on the first
line and $f$ with $e\cdot f=1$, then replace $f$ by
$f-(f\cdot f)e/2$. In the resulting basis the norm is $2xy$,
so the isotropic lines are precisely the two coordinate axes.
Counting hyperplanes of $C/G$ gives the stated neighbor count.

Each normalized sign vector outside $C$ but in $G^\perp$ is isotropic,
since its norm is $N=0$ modulo $p$. It belongs to exactly one of
these neighbors, namely
$(C\cap v^\perp)+\langle v\rangle$. Thus the total number of
new sign classes, counted over the neighbors, is $M_G-r_C$.
An old sign class in $G$ is never lost. An old class outside $G$
is excluded by precisely $p^{h-1}$ hyperplanes containing $G$,
the difference between the hyperplane counts in dimensions
$h$ and $h-1$. Summing gains minus losses proves the identity.
\end{proof}

This is a finite incidence identity in the classical geometry of
self-dual code neighbors. It is not a general positivity theorem.
In particular, at $G=0$ its right-hand side is
$2^{N-1}-r_C(1+p^{m-1})$, which can be negative even when the
code has far fewer than $N$ sign classes. Unrestricted averaging
therefore does not guarantee a useful move.

\begin{corollary}[A modular pool splits into two codes]
\label{code:two-codes}
If $G$ has dimension $m-1$, there is exactly one other self-dual
code $C'$ containing $G$, and
\[
 \{ \text{sign classes in }G^\perp\}=S(C)\cup S(C'),\qquad
 S(C)\cap S(C')=S(G),\qquad
 |S(C')|=M_G-r_C+b_G.
\]
If $G$ contains five orthogonal sign rows, the sign words of each
code separately form a real orthogonal frame.
\end{corollary}
\begin{proof}
An isotropic residue in $G^\perp$ projects to one of the two
isotropic lines of $G^\perp/G$, hence belongs to $C$ or $C'$.
Their intersection is $G$. Count the union, and apply
Corollary~\ref{code:five} for real orthogonality.
\end{proof}

This removes a compatibility search in the codimension-one case.
For the terminal twelve-row order-twenty certificate, the ten rows
retained after deleting indices $3,8$ span such a $G$, and exact
enumeration gives
\[
 r_C=12,\qquad b_G=10,\qquad M_G=22,\qquad |S(C')|=20.
\]
The twelve exact extensions of the retained frame consist of the two
deleted old classes and the ten new classes. The two-code theorem,
rather than an independent guess that all candidates are compatible,
explains the resulting completion.

The checker also verifies the incidence identity for all four
neighbors containing a specified four-dimensional subspace of a
self-dual ternary code at order twelve. These are finite examples.
An all-order lower bound on the modular pool, or a guarantee that
some choice of retained space makes the displayed surplus positive,
has not been established.
\section{Modular lifting in the orders congruent to four modulo eight}
\label{sec:four-mod-eight-lifting}

The preceding lifting theorem used $N=4(q+1)$ with odd $q$, so its
orders were divisible by eight. The complementary choice $N=4q$
gives an especially simple result throughout the other admissible
congruence class. Its only spurious modular extensions, once five
orthogonal rows are present, are the rows already installed.
This is a counting and lifting statement; it does not guarantee the
surplus modular solutions needed for a new row. No novelty claim is
made pending a targeted prior-art check.

Let $P$ have $r$ orthogonal sign rows $p_i$ of length $N$. Count sign
vectors up to negation, and write $M_q(P)$ for the number satisfying
$Pv\equiv0\pmod q$ and $E(P)$ for the number satisfying $Pv=0$.

\begin{theorem}[The row correction at orders $4\bmod8$]
Let $q\ge3$ be any odd integer, $N=4q$, and $r\ge5$. Then
\[
 \boxed{E(P)=M_q(P)-r.}
\]
Every modular sign extension which is not exact is the sign or
negative of exactly one existing row. Primality of $q$ is not required.
\end{theorem}
\begin{proof}
For a modular sign vector $v$ put $d_i=p_i\cdot v$.
Orthogonal rows differ in $N/2=2q$ positions, an even number, so
all $d_i$ have the same residue modulo four. Every $d_i$ is even
and divisible by odd $q$.

If that common residue is two, each $d_i$ must be $2q$ or $-2q$,
since the next possible magnitude $6q$ exceeds $N$. But Bessel gives
\[
 4q^2r=\sum_i d_i^2\le N\|v\|^2=N^2=16q^2,
\]
contradicting $r\ge5$. Hence all $d_i$ are zero modulo four and
therefore belong to $\{0,N,-N\}$. A nonzero value means
$v=\pm p_i$, by equality in the sign-vector inner product bound.
All other inner products then vanish by row orthogonality.
Conversely every existing row is a modular sign vector because its
nonzero self-inner-product is $N=4q$. The $r$ row classes are distinct,
proving the count formula.
\end{proof}

\subsection{The sharp four-row boundary and closed quadruples}
The five-row threshold has an exact geometric explanation. Call four
orthogonal sign rows a \emph{closed quadruple} if their entrywise
product is constant. This is standard terminology in Hadamard
switching; see Orrick, \emph{Switching Operations for Hadamard
Matrices}, Section 3.1 and Proposition 3.1,
\url{https://arxiv.org/abs/math/0507515}.
In particular the prohibition of such a quadruple inside a full
Hadamard matrix of order $4\bmod8$ is classical.

\begin{proposition}[Exact boundary correction]
For four orthogonal sign rows of length $N=4q$, $q\ge3$ odd,
\[
 E(P)=M_q(P)-4-4\,\mathbf1_{\{P\text{ is closed}\}}.
\]
When the quadruple is closed, the four extra modular classes are
represented by half-sums
\[
 v=\frac12\sum_{i=1}^4\sigma_i p_i,
 \qquad\sigma_i\in\{\pm1\},\qquad
 \prod_i\sigma_i=-c,
\]
where $c$ is the constant entrywise row product. Negating all four
$\sigma_i$ gives the same antipodal class.
\end{proposition}
\begin{proof}
The common-zero residue case in the preceding proof still produces
only exact extensions and existing rows. In the common-two case,
all four inner products are $\sigma_iN/2$, and Bessel is saturated.
Thus $v$ has exactly the stated half-sum representation, with no
orthogonal-complement component. Four signs sum to $+2$ or $-2$
exactly when their product is $-1$. Consequently every coordinate
of this half-sum is a sign if and only if
$\prod_i p_{ij}=-\prod_i\sigma_i$ for every coordinate $j$.
This is exactly the closed condition with the displayed restriction.
Eight choices of $\sigma$ modulo simultaneous negation give four
classes; independence of the four rows makes their half-sums distinct.
\end{proof}

For example $H_4\otimes\mathbf1_q$ is a closed quadruple with no
exact extension, yet has eight modular classes: four installed rows
and four half-sums. Thus subtracting only the original rows would
give a false positive at this sharp boundary. The no-extension claim
follows directly from its four length-$q$ block sums: orthogonality
forces all sums zero, impossible for odd block size.

\subsection{Finite algebra and the remaining positivity problem}
For any integer modulus $q$, use the integral group algebra of
$(\mathbb Z/q\mathbb Z)^r$ and let $c_j$ be column $j$ of $P$.
Then
\[
 C_q(P)=[X^0]\prod_{j=1}^N(X^{c_j}+X^{-c_j})=2M_q(P).
\]
Under the five-row hypothesis,
\[
 \boxed{E(P)=\tfrac12C_q(P)-r.}
\]
Hence $C_q(P)>2r$ is an exact sufficient condition for a new sign row,
and $C_q(P)=2r$ certifies terminality. The group-algebra formula
works for composite odd $q$ as well as primes. For prime $q$ it can
also be computed by finite-field Gaussian elimination or a Fourier
character sum. When composite moduli are used, ordinary field row
reduction must be replaced by valid arithmetic over the residue ring
or a suitable Chinese-remainder procedure.

For prime $q$, the row code $C\subseteq\mathbb F_q^N$ generated by
$P$ is self-orthogonal, because $PP^{\mathsf T}=NI_r\equiv0\pmod q$.
If its rank is $s$, then $s\le N/2$. An information-coordinate search
therefore uses $2^{N-s-1}$ sign assignments; large real row count
does not imply equally large modular rank. The existing rows already
give $M_q(P)\ge r$. Proving the strict surplus needed for extension
requires additional code structure, rather than rank alone.

Exact examples from explicit Hadamard matrices illustrate the formula:
\[
\begin{array}{c|r|r|r|r}
N&r&\operatorname{rank}_{q}P&M_q(P)&E(P)\\\hline
12&5&5&18&13\\
12&6&6&12&6\\
12&10&6&12&2\\
20&10&10&20&10\\
20&18&10&20&2\\
28&14&14&28&14\\
28&26&14&28&2
\end{array}
\]
Here $q=N/4$ and the matrices are constructed explicitly in the
standard-library checker. An independent integer group-algebra
calculation gives $C_3(P)=36$ for the five-row order-twelve example,
so $36/2-5=13$ exactly. Its thirteen candidate classes need not all
be mutually orthogonal: counting single-row extensions alone does
not automatically supply a complete frame.

Together, the $N=4(q+1)$ theorem for odd $q$ and the present
$N=4q$ theorem treat both admissible residue classes modulo eight,
with different explicit corrections. Neither asserts the requisite
modular positivity uniformly, so neither proves the full Hadamard
conjecture. Their concrete value is that they replace an approximate
modular extension problem by an exact real one after accounting for
fully classified exceptional configurations.

\section{Concentration moves and target-dependent bridges}

Let $\mathcal P_d$ be the probability laws on the sign cube with
zero first and second moments. Write $A$ for the matrix evaluating
the constant, degree-one and degree-two characters, and $b=(1,0,\ldots,0)$.
Thus $\mathcal P_d=\{p\ge0:Ap=b\}$ and $C(p)=\|p\|_2^2$.

\begin{theorem}[An exact monotone support-reduction algorithm]
Given any rational feasible law $p$ with support $S$, repeat the
following operation while the columns $A_S$ are dependent.
Choose nonzero $z\in\ker A_S$, reverse its sign if necessary so
$p\cdot z\ge0$, and set
\[
 t=\min_{j:z_j<0}\frac{p_j}{-z_j},\qquad p\leftarrow p+tz.
\]
Every intermediate law remains nonnegative and exactly feasible,
collision increases strictly at each nontrivial step, and the entire
interpolating path has nondecreasing collision. After at most
$|S|-\operatorname{rank}A_S$ steps it reaches a vertex supported on
at most $1+d+\binom d2$ atoms. All calculations can be rational.
\end{theorem}
\begin{proof}
The constant row of $Az=0$ gives $\sum z_j=0$, so $z$ has both
signs and $t>0$ is finite. The definition of $t$ preserves positivity
and kills at least one atom; moment feasibility follows from $Az=0$.
For $0\le u\le t$,
$C(p+uz)=C(p)+2u(p\cdot z)+u^2\|z\|^2$ is nondecreasing and
strictly increases at the endpoint. The feasible face containing
$p$ in its relative interior has dimension
$|S|-\operatorname{rank}A_S$. Each boundary step passes to a proper
face, decreasing that dimension at least one. At a zero-dimensional
face the active columns are independent. This proves termination
and the support bound. One can choose a minimally supported kernel
vector, making every step a circuit move. This is the elementary
convex-objective endpoint principle familiar from pipage-type
rounding~\cite{AS}, not a claimed new general optimization method.
\end{proof}
The algorithm can terminate at a suboptimal vertex, including the
weighted four-bit examples. It does not guarantee minimal entropy
or a Hadamard construction. The result identifies exactly what local
support deletion can do without a collision barrier.

\begin{theorem}[Sharp barrier for a specified straight bridge]
Let $p,q\in\mathcal P_d$ be distinct and set
$a=C(p)$, $b=C(q)\ge a$, $c=p\cdot q$, $e=a+b-2c>0$.
The feasible bridge $p_t=(1-t)p+tq$ has minimum collision
\[
 \min_{0\le t\le1}C(p_t)
 =a-\frac{(a-c)_+^2}{a+b-2c}.
\]
The minimizer is $t=(a-c)_+/e\le1/2$. Its collision loss is at most
$a^2/(a+b)\le a/2$, with the first bound attained for disjoint
supports. If $c\ge a$ the bridge is monotone and has no loss.
\end{theorem}
\begin{proof}
Expand $C(p_t)=a+2(c-a)t+et^2$. If $c<a$, the minimizer is
$(a-c)/e$, at most one half since $b\ge a$. Since $c\ge0$ and
$c<a\le b$, the function $(a-c)^2/(a+b-2c)$ decreases with $c$:
its derivative is $-2(a-c)(b-c)/(a+b-2c)^2\le0$.
Its value at zero is $a^2/(a+b)$. If $c\ge a$ the quadratic is
nondecreasing on the entire interval.
\end{proof}
If an explicit optimal law $q$ is available, this gives an actual
nonlocal feasible route and a quantitative loss budget. It does not
provide the missing optimal target. Translating $q$ by cube signs
preserves feasibility and collision, and averaging those translates
gives the uniform cube. Consequently some translate has overlap
$p\cdot q\ge2^{-d}$, which may improve the displayed bridge bound.

\section*{A sharp bridge need not be a global barrier}
For $d=5$, use the eight-word linear Walsh code
\[
 C=\{((-1)^{u\cdot m})_{m=1}^5:u\in\mathbb F_2^3\}.
\]
The masks $1,\ldots,5$ include a basis and are distinct and nonzero,
so the code is injective and every coset's uniform law is unbiased
and pairwise independent. There are four disjoint cosets in the
32-word cube. Let $q_1,q_2,q_3$ be three of these uniform laws and
$p=(q_1+q_2)/2$. Then
\[
 C(p)=1/16,\quad C(q_3)=1/8,\quad p\cdot q_3=0.
\]
The straight bridge to $q_3$ has minimum $1/24$ at $t=1/3$ and
collision barrier $1/48$, attaining the sharp general bound.
Nevertheless the bridge from $p$ to $q_1$ has overlap $p\cdot q_1=1/16$
and collision $1/16+t^2/16$, which is monotone. Both targets have
optimal collision $1/8$ by the residue-sensitive bound.
Thus even an exactly sharp barrier to a selected optimum cannot be
promoted to a barrier for reaching \emph{any} optimum. The globally
optimal four-bit barrier required its additional section argument;
there is no automatic extension to larger moment polytopes.

\begin{thebibliography}{9}
\bibitem{AS} A. A. Ageev and M. I. Sviridenko, Pipage rounding:
a new method of constructing algorithms with proven performance
guarantee, Journal of Combinatorial Optimization 8 (2004), 307--328.
\url{https://doi.org/10.1023/B:JOCO.0000038913.96607.c2}.
\end{thebibliography}
\section{Quantitative low-degree barriers and support repair}
\label{sec:nonlinear-quantitative}
The statements concern uniform laws on $N$ binary seed labels, with
$N\equiv4\pmod8$, $r\ge4$, and Boolean output degree at most $r-2$.
They do not automatically constrain nonuniform measures or larger seed
models. Put $S\subset\mathbb F_2^r$, $|S|=N$, $T=S^c\ne\varnothing$,
and $a=r-\operatorname{affdim}(T)$. The preceding pass proved that the
even restriction code of $\mathrm{RM}(r-2,r)$ has binary bilinear rank
$N-2a-2$. We recall the reason: its dual consists of affine functions
supported on $S$, of dimension $a$; the even code radical is that dual
plus the constant function. Reed--Muller duality and the Boolean-code
background can be found in Carlet's primary monograph,
\url{https://www.math.univ-paris13.fr/~carlet/book-fcts-Bool-vect-crypt-codes.pdf}.
All new estimates below have explicit proofs; priority is unestablished.

\begin{theorem}[A quantitative moment-error floor]
\label{thm:nonlinear-moment-floor}
Suppose $m=N-4$ Boolean sign observables $X_1,\ldots,X_m$ have exact
zero means on $S$, with degree at most $r-2$. They need not be pairwise
independent. At least $\max\{0,2a-3\}$ unordered pairs have nonzero
correlation. In particular their augmented moment error satisfies
\[
 E:=\left\|\mathbb E[(1,X)(1,X)^T]-I_{m+1}\right\|_F^2
 \ge\frac{32}{N^2}\max\{0,2a-3\}.
\]
For four omitted labels and $r\ge5$, this gives
$E\ge32(2r-9)/N^2$, since $a\ge r-3$.
\end{theorem}
\begin{proof}
Write $B_{zi}=(1-X_i(z))/2$ and $G=B^TB$ over $\mathbb F_2$.
Balanced columns have weight $N/2$, even, and belong to the even
restriction code, so $\operatorname{rank}G\le N-2a-2$.
The exact pairwise-independent target would be $G_*=I+J_m$,
nonsingular because $m$ is even. Let $q_{ij}=\sum_{z\in S}X_i(z)X_j(z)$.
Balancedness gives $q_{ij}\in4\mathbb Z$ and
\[
 G_{ij}=1+q_{ij}/4\pmod2\quad(i\ne j),\qquad G_{ii}=0.
\]
Thus $G=I+J+A_H$, where $H$ is the graph of pairs with $q_{ij}/4$ odd.
Let $e$ be its number of edges. Each nontrivial connected component of
$H$ with $v$ vertices has at least $v-1$ edges, while its $I+A_H$ block
has rank at least one. Isolated vertices contribute no nullity. Hence
$\operatorname{nullity}(I+A_H)\le e$. Adding the rank-one matrix $J$
gives $\operatorname{nullity}G\le e+1$.
But the rank bound gives $\operatorname{nullity}G\ge2a-2$, so
$e\ge2a-3$. Each such pair has $|q_{ij}|\ge4$; the Frobenius error
counts every unordered correlation twice. This proves the bound.
\end{proof}

The edge-count estimate is sharp at the binary-Gram level: if
$q=2a-3$ is positive and $m\ge2q$, take $H$ to be a matching of $q$
edges. Since $q$ is odd, $I+J+A_H$ has nullity $q+1=2a-2$.
Indeed its kernel vectors have common unmatched-coordinate value $u$,
and the sum on each matched pair equals $u$; the $q$ pair choices and
$u$ are free. This does not assert that the matching Gram is realized
by balanced low-degree truth columns with the smallest possible real
correlations. Thus no sharpness claim is made for the moment-error
constant as a distribution theorem.

\begin{theorem}[Minimum support edits to unlock degree repair]
\label{thm:nonlinear-minimum-edits}
For this support operation no congruence condition on $N$ is needed.
Assume $|T|\ge r+1$. The smallest number of one-for-one seed support
exchanges needed to make the omitted labels affinely span the entire
$r$-space is exactly $a=r-\operatorname{affdim}(T)$.
For any already-known exact law carried on $S$, its output truth rows
can be transferred through these exchanges while preserving its entire
observable distribution, hence every moment and its collision probability.
All intermediate truth functions can have degree at most $r-1$; at the
final support they can have degree at most $r-2$.
\end{theorem}
\begin{proof}
One replacement raises affine dimension by at most one, giving the
lower bound $a$. For the upper bound, retain an affine basis of $T$,
of size $r-a+1$. There are at least $a$ redundant points because
$|T|\ge r+1$. Remove redundant points one at a time and replace them
with points outside the current affine span. Each replacement raises
the span dimension by one; after $a$ exchanges it is full.

A replacement of omitted point $t$ by included point $s$ replaces
included seed label $s$ by $t$. Transfer the output row at $s$ to $t$,
leaving every other output row fixed. This changes no observable atom
or weight. Extend each Boolean truth function by zero on omitted labels
and change its value at one omitted point to make total parity zero;
this gives degree at most $r-1$. At the final support, choose an affine
basis in the complement and change only those outside values to cancel
the $r+1$ affine parity syndromes. The corresponding $(r+1)$-square
matrix of affine evaluations is invertible, so the repair exists and
lies in $\mathrm{RM}(r-2,r)$. No truth value at a positive-mass seed is
altered by these outside repairs.
\end{proof}
This operation escapes a low-rank \emph{seed encoding}. It preserves the
observable law rather than producing a new or more concentrated law,
so it is not a missing-distribution construction.

\subsection{Four omitted labels: Fourier sparsity versus capacity}
Let $N=2^r-4$ and let $\mu$ be uniform on $S$. For a nonzero character
$c$, the full-cube sum vanishes and
\[
 \widehat\mu(c)=\mathbb E_\mu(-1)^{c\cdot z}
 =-\frac1N\sum_{z\in T}(-1)^{c\cdot z}.
\]
Four distinct omitted points have two affine types. Exact character
counting gives:
\[
 \begin{array}{c|c|c|c}
 T\text{ type}&\operatorname{affdim}T&
 \#\{c\ne0:\widehat\mu(c)\ne0\}&
 \text{even restriction bilinear rank}\\\hline
 \text{affine plane}&2&2^{r-2}-1&N-2r+2\\
 \text{affine-independent tetrahedron}&3&5\cdot2^{r-3}-1&N-2r+4
 \end{array}
\]
For the plane translate to $\{0,e_1,e_2,e_1+e_2\}$; its character sum
is four precisely when the first two frequency coordinates vanish.
For the tetrahedron use $\{0,e_1,e_2,e_3\}$; its sum is four, two,
zero, or minus two according as the first three frequency coordinates
have weights zero, one, two, or three. Translation only changes signs.
Thus the Fourier-sparser seed is also the lower-capacity low-degree
encoding. At twelve labels the two cases have three versus nine nonzero
seed Fourier modes and, by the preceding exact classification, five
versus seven possible quadratic unbiased pairwise-independent bits.
Neither reaches the eight-bit extraction target.

The table is an exact narrow-model tradeoff, not a general uncertainty
bound or an argument that many Fourier modes guarantee attainment.
For $|T|=4<r+1$, the support-edit theorem cannot unlock a degree-$r-2$
full basis at all: no four omitted labels can span the seed space when
$r\ge4$. Higher output degree, a different seed size, or a different
weight model is necessary.

\section{The unresolved step}
The modular correction is now explicit for both order classes modulo
eight, but its positivity is not established in general.
A self-orthogonal code makes its sign words compatible once five
orthogonal rows are installed, but does not force many such words.
Self-dual neighbor moves solve a concrete artificial restriction at
order twenty; the general rule still requires a real extension and
cannot preserve all rows of a terminal frame while augmenting it.
The concentration endpoint algorithm can terminate at a suboptimal
vertex, and support relabeling preserves rather than constructs a law.

The central missing ingredient is therefore an actual construction
or lower-count theorem, or an exchange theorem that can replace
incompatible rows globally. No universal extension assertion is
assumed here. Finite checks validate explicit benchmarks; they do
not prove unknown-order existence or establish priority.
\end{document}
